% Semaphore UI文档
% Semaphore_UI.tex

\documentclass[12pt,UTF8]{ctexbook}

% 设置纸张信息。
\input{../../../../../../common/page}

% 设置字体，并解决显示难检字问题。
\xeCJKsetup{AutoFallBack=true}
% 注意：ctexbook类已默认设置SimSun为CJKrmdefault
% 以下设置用于确保字体回退和扩展字体可用
% 仅设置扩展字体以避免与默认设置冲突
\setCJKfamilyfont{hei}{SimHei}
\setCJKfamilyfont{kai}{KaiTi}

% 目录 chapter 级别加点（.）。
\usepackage{titletoc}
\titlecontents{chapter}[0pt]{\vspace{3mm}\bf\addvspace{2pt}\filright}{\contentspush{\thecontentslabel\hspace{0.8em}}}{}{\titlerule*[8pt]{.}\contentspage}

% 支持目录点击跳转
\usepackage[colorlinks,linkcolor=blue,citecolor=blue,urlcolor=blue]{hyperref}

% 设置 part 和 chapter 标题格式。
\ctexset{
	part/name= {第,卷},
	part/number={\chinese{part}},
	chapter/name={第,篇},
	chapter/number={\arabic{chapter}}
}

% 图片相关设置。
\usepackage{graphicx}
\graphicspath{{Images/}}

% 设置署名格式。
\newenvironment{shuming}{\hfill\zihao{4}}

% 注脚每页重新编号，避免编号过大。
\usepackage[perpage]{footmisc}

% 设置古文原文格式。
\newenvironment{yuanwen}{\bfseries\zihao{4}}

% 设置段落间距
\setlength{\parskip}{0.8em plus 0.2em minus 0.1em}

% 列表格式支持，列表项向右偏移。
\usepackage{enumitem}

\begin{document}

\frontmatter
\chapter{前言、序言}

\mainmatter

\chapter{Semaphore UI概述}

\section{什么是Semaphore UI}

\textbf{英文原文：}

Semaphore UI is a modern, open-source web UI and API for running automation with Ansible, Terraform/OpenTofu, PowerShell, Shell/Bash, and Python. 

Semaphore is written in Go (lightweight, fast) and runs on Windows, macOS, and Linux. It supports SQLite, MySQL, and PostgreSQL.

\textbf{中文翻译：}

Semaphore UI是一个现代化的、开源的Web用户界面和API，用于运行Ansible、Terraform/OpenTofu、PowerShell、Shell/Bash和Python的自动化任务。

Semaphore使用Go语言编写（轻量级、快速），可在Windows、macOS和Linux系统上运行。它支持SQLite、MySQL和PostgreSQL数据库。

\section{首次运行检查清单（从零到第一个任务）}

\textbf{英文原文：}

First run checklist (from zero to first task) 

Use this as a simple "happy path" to get productive quickly: 

\begin{enumerate}[leftmargin=3.5em]
    \item Install Semaphore using your preferred method: Installation 
    \item Run initial setup (config + first admin user): Interactive setup 
    \item Create a project to isolate work (teams, infra, apps): Projects 
    \item Connect what Semaphore needs 
    \begin{itemize}[leftmargin=2em]
        \item Source: Repositories 
        \item Secrets/credentials: Key Store 
        \item Targets: Inventory 
        \item Variables: Variable Groups 
    \end{itemize}
    \item Create a task template and run it 
    \begin{itemize}[leftmargin=2em]
        \item App-specific guides: Ansible, Terraform/OpenTofu, Shell/Bash 
        \item Run and monitor: Tasks 
    \end{itemize}
    \item Automate \& operationalize 
    \begin{itemize}[leftmargin=2em]
        \item Schedule runs: Schedules 
        \item Control access: RBAC 
        \item Get alerts: Notifications
    \end{itemize}
\end{enumerate}

\textbf{中文翻译：}

首次运行检查清单（从零到第一个任务）

使用这个简单的"快乐路径"快速上手：

\begin{enumerate}[leftmargin=3.5em]
    \item 使用您偏好的方法安装Semaphore：安装指南
    \item 运行初始设置（配置 + 第一个管理员用户）：交互式设置
    \item 创建项目来隔离工作（团队、基础设施、应用）：项目管理
    \item 连接Semaphore所需的内容
    \begin{itemize}[leftmargin=2em]
        \item 源代码：代码仓库
        \item 密钥/凭据：密钥存储
        \item 目标：清单管理
        \item 变量：变量组
    \end{itemize}
    \item 创建任务模板并运行它
    \begin{itemize}[leftmargin=2em]
        \item 应用特定指南：Ansible、Terraform/OpenTofu、Shell/Bash
        \item 运行和监控：任务管理
    \end{itemize}
    \item 自动化与运营化
    \begin{itemize}[leftmargin=2em]
        \item 调度运行：计划任务
        \item 控制访问：基于角色的访问控制（RBAC）
        \item 获取警报：通知功能
    \end{itemize}
\end{enumerate}

\section{关键概念（术语表）}

\textbf{英文原文：}

Key concepts (glossary) 

If you're new, these terms show up everywhere in the UI: 

\begin{itemize}[leftmargin=3.5em]
    \item \textbf{Project}: the main unit of separation (teams/infrastructures/applications) — Projects 
    \item \textbf{Repository}: where your playbooks/modules/scripts live — Repositories 
    \item \textbf{Inventory}: hosts/groups/connection settings for Ansible-style runs — Inventory 
    \item \textbf{Variable Group (Environment)}: reusable variables and configuration per project — Variable Groups 
    \item \textbf{Key Store}: encrypted credentials (SSH keys, tokens, passwords) — Key Store 
    \item \textbf{Task / Task template}: the definition and the run of automation — Tasks 
    \item \textbf{Runner}: where tasks execute (local or remote) — Runners
\end{itemize}

\textbf{中文翻译：}

关键概念（术语表）

如果您是新手，这些术语在用户界面中随处可见：

\begin{itemize}[leftmargin=3.5em]
    \item \textbf{项目（Project）}：主要的隔离单元（团队/基础设施/应用）— 项目管理
    \item \textbf{仓库（Repository）}：存放您的playbooks/模块/脚本的地方 — 仓库管理
    \item \textbf{清单（Inventory）}：用于Ansible风格运行的主机/组/连接设置 — 清单管理
    \item \textbf{变量组（环境）}：每个项目可重用的变量和配置 — 变量组管理
    \item \textbf{密钥存储（Key Store）}：加密的凭据（SSH密钥、令牌、密码）— 密钥存储管理
    \item \textbf{任务/任务模板}：自动化的定义和运行 — 任务管理
    \item \textbf{运行器（Runner）}：任务执行的地方（本地或远程）— 运行器管理
\end{itemize}

\section{常见工作流程}

\textbf{英文原文：}

Common workflows 

\begin{itemize}[leftmargin=3.5em]
    \item \textbf{CI/CD in Semaphore}: build, deploy, and rollback — CI/CD 
    \item \textbf{Run at scale}: distribute execution using runners — Runners 
    \item \textbf{Programmatic automation}: integrate Semaphore into your pipelines — API 
    \item \textbf{Manage via CLI}: setup, users, runners, migrations — CLI
\end{itemize}

\textbf{中文翻译：}

常见工作流程

\begin{itemize}[leftmargin=3.5em]
    \item \textbf{Semaphore中的CI/CD}：构建、部署和回滚 — CI/CD流程
    \item \textbf{大规模运行}：使用运行器分发执行 — 运行器管理
    \item \textbf{编程式自动化}：将Semaphore集成到您的流水线中 — API集成
    \item \textbf{通过CLI管理}：设置、用户、运行器、迁移 — 命令行界面
\end{itemize}

\chapter{安装指南}

\section{安装方式概述}

\textbf{英文原文：}

Installation 

You can install Semaphore in multiple ways, depending on your operating system, environment, and preferences: 

\begin{itemize}[leftmargin=3.5em]
    \item \textbf{Package manager} 
    Install Semaphore using a native package for your distribution (e.g., apt for Debian/Ubuntu or dnf for RHEL-based systems). This is the easiest way to get started on Linux servers and integrates well with system services. 
    Learn more » 
    
    \textbf{Additional notes:}
    
    \begin{itemize}[leftmargin=2em]
        \item \textbf{Tip}: Look into the manual installation on how to set-up your Python/Ansible/Systemd environment! 
        \item \textbf{Download package file}: Download package file from Releases page. 
        \item \textbf{Package types}: 
        \begin{itemize}[leftmargin=1.5em]
            \item Debian / Ubuntu (x64)
            \item Debian / Ubuntu (ARM64)
            \item CentOS (x64)
            \item CentOS (ARM64)
        \end{itemize} 
        \item \textbf{Installation commands}: Here are several installation commands, depending on the package manager.
        
        \textbf{Debian / Ubuntu (x64)}:
        \begin{verbatim}
        wget https://github.com/semaphoreui/semaphore/releases/download/v2.17.15/semaphore_2.17.15_linux_amd64.deb
        
        sudo dpkg -i semaphore_2.17.15_linux_amd64.deb
        \end{verbatim}
        
        \textbf{Setup and Run Semaphore}:
        \begin{verbatim}
        semaphore setup
        
        semaphore server --config=./config.json
        \end{verbatim}
        
        Semaphore will be available via this URL https://localhost:3000.
    \end{itemize}
    
    \item \textbf{Docker} 
    Run Semaphore as a container using Docker or Docker Compose. Ideal for fast setup, sandboxed environments, and CI/CD pipelines. Recommended for users who prefer infrastructure as code. 
    Learn more » 
    
    \item \textbf{Cloud} 
    Guidance for deploying Semaphore to cloud platforms using VMs, containers, or Kubernetes with managed services. 
    Learn more » 
    
    \item \textbf{Binary file} 
    Download a precompiled binary from the releases page. Great for manual installation or embedding in custom workflows. Works across Linux, macOS, and Windows (via WSL). 
    Learn more » 
    
    \item \textbf{Kubernetes (Helm chart)} 
    Deploy Semaphore into a Kubernetes cluster using Helm. Best suited for production-grade, scalable infrastructure. Supports easy configuration and upgrades via Helm values. 
    Learn more » 
\end{itemize}

See also: 

\begin{itemize}[leftmargin=3.5em]
    \item Run as service 
    \item Manual installation 
\end{itemize}

\textbf{中文翻译：}

安装指南

您可以根据操作系统、环境和偏好以多种方式安装Semaphore：

\begin{itemize}[leftmargin=3.5em]
    \item \textbf{包管理器} 
    使用适用于您的发行版的本地包安装Semaphore（例如，Debian/Ubuntu使用apt，基于RHEL的系统使用dnf）。这是在Linux服务器上开始使用的最简单方法，并且与系统服务集成良好。
    了解更多 » 
    
    \textbf{补充说明：}
    
    \begin{itemize}[leftmargin=2em]
        \item \textbf{提示}：请查阅手动安装指南以了解如何设置Python/Ansible/Systemd环境！
        \item \textbf{下载包文件}：从发布页面下载包文件。
        \item \textbf{包类型}：Debian和Ubuntu使用*.deb，CentOS和RedHat使用*.rpm。
        \item \textbf{安装命令}：根据包管理器不同，有几种安装命令。
        
        \textbf{Debian / Ubuntu (x64)}:
        \begin{verbatim}
        wget https://github.com/semaphoreui/semaphore/releases/download/v2.17.15/semaphore_2.17.15_linux_amd64.deb
        
        sudo dpkg -i semaphore_2.17.15_linux_amd64.deb
        \end{verbatim}
        
        \textbf{设置和运行Semaphore}:
        \begin{verbatim}
        semaphore setup
        
        semaphore server --config=./config.json
        \end{verbatim}
        
        Semaphore将通过以下URL访问：https://localhost:3000。
    \end{itemize}
    
    \item \textbf{Docker} 
    使用Docker或Docker Compose将Semaphore作为容器运行。适用于快速设置、沙盒环境和CI/CD流水线。推荐给偏好基础设施即代码的用户。
    了解更多 » 
    
    \item \textbf{云平台} 
    使用虚拟机、容器或Kubernetes与托管服务将Semaphore部署到云平台的指南。
    了解更多 » 
    
    \item \textbf{二进制文件} 
    从发布页面下载预编译的二进制文件。适用于手动安装或嵌入到自定义工作流中。可在Linux、macOS和Windows（通过WSL）上运行。
    了解更多 » 
    
    \item \textbf{Kubernetes（Helm图表）} 
    使用Helm将Semaphore部署到Kubernetes集群中。最适合生产级、可扩展的基础设施。支持通过Helm值进行简单配置和升级。
    了解更多 » 
\end{itemize}

另请参阅：

\begin{itemize}[leftmargin=3.5em]
    \item 作为服务运行
    \item 手动安装
\end{itemize}

\section{安装额外的Python包}

\textbf{英文原文：}

Installing Additional Python Packages 

Some Ansible modules and roles require additional python packages to run. To install additional python packages, create a requirements.txt file and mount it in the /etc/semaphore directory on the container. For example, you could add the following lines to your docker-compose.yml file: 

volumes: 
  - /path/to/requirements.txt:/etc/semaphore/requirements.txt 

The packages specified in the requirements file will be installed when the container starts up. 

For more information about Python requirements files, see the Pip Requirements File Format reference

\textbf{中文翻译：}

安装额外的Python包

一些Ansible模块和角色需要额外的Python包才能运行。要安装额外的Python包，请创建一个requirements.txt文件并将其挂载到容器中的/etc/semaphore目录。例如，您可以将以下行添加到docker-compose.yml文件中：

volumes: 
  - /path/to/requirements.txt:/etc/semaphore/requirements.txt 

requirements文件中指定的包将在容器启动时安装。

有关Python需求文件的更多信息，请参阅Pip需求文件格式参考。

\end{document}